\section{Complete short-channel catalogues and the full period group}
\label{sec:canonical}
The finite-index theorem concerns an arbitrarily selected collection.
Its ambiguity need not survive when all short visible normal channels
are recorded. We first prove a geometric completeness principle, valid
without analyticity or asymmetry, and then combine it with the local
inverse. The resulting global theorem requires no primitive cycle pair.
The word complete is a condition on the acquisition protocol, not a
statistical conclusion drawn from a list with unknown omissions.

\subsection{The lifted graph of shortest clear bridges}
Write $\mathcal C=\{C_i+\lambda:1\le i\le r,\lambda\in\Lambda\}$
for all physical obstacle copies and $\mathcal O=\bigcup_{C\in\mathcal C}C$.
Its covering radius is
\begin{equation}\label{eq:covering-radius}
 \rho_* =\sup_{x\in\R^2}\dist(x,\mathcal O)<\infty.
\end{equation}
Finiteness follows from periodicity and the presence of an obstacle.
It bounds distance to a reflecting body, not the length of every free
flight: infinite-horizon tables are included.

For distinct copies $C,D$ let $\ell(C,D)=\dist(C,D)>0$.
Strict convexity makes their shortest connecting segment unique.
Indeed, the closest difference vector is unique by strict convexity of
the squared Euclidean norm, and the two support points in its normal
direction are unique. Denote this segment by $[p_{CD},q_{CD}]$.
It is normal to both boundaries. Call it clear if its open segment meets
no other obstacle, including at a tangency. A clear segment has positive
clearance from all third bodies, since the family is locally finite.
The endpoints also have positive separation from third bodies.

For $R>0$ let $\mathcal G_R$ have vertex set $\mathcal C$ and an
undirected edge between $C,D$ precisely when $\ell(C,D)<R$ and their
shortest segment is clear. Its quotient by translations in $\Lambda$
is finite; loops and multiple edges in the quotient are retained.
There are only finitely many pair types with $\ell(C,D)<R$, because
a bounded bridge and the bounded diameters bound their relative period
translation. In particular, the set of their distinct gap lengths is
finite even though the lifted graph has infinitely many vertices.

\begin{lemma}[Obstruction descent]\label{lem:obstruction-descent}
For each $R$, the graph $\mathcal G_R$ has the same connected components
as the graph that joins all pairs with gap less than $R$, whether or
not their shortest segments are clear.
\end{lemma}
\begin{proof}
Order the finite set of gap lengths less than $R$. If a shortest segment
from $C$ to $D$ meets a third body $F$ at $z$, disjointness implies
$z$ is strictly between its endpoints. Consequently
\[
 \ell(C,F)\le |p_{CD}-z|<\ell(C,D),\qquad
 \ell(F,D)\le |z-q_{CD}|<\ell(C,D).
\]
This also holds when the intersection is a single tangency point.
Induction over the ordered gap lengths replaces the two shorter pairs
by clear-edge paths. Concatenating them replaces the original pair by
a path in $\mathcal G_R$. The induction starts at the smallest gap,
which cannot be obstructed. No termination assumption on an infinite
sequence of distinct lengths is needed: periodicity has made the set
of relevant lengths finite. Every clear edge is also an edge of the
unrestricted gap graph, proving the converse inclusion of components.
\end{proof}

\begin{theorem}[Geometric completeness at bounded range]
\label{thm:geometric-completeness}
For every periodic family of disjoint strictly convex bodies and every
$R>2\rho_*$, the lifted graph $\mathcal G_R$ is connected. Its finite
quotient visits every obstacle orbit, is connected, and its cycle
translations generate the entire period lattice over $\Z$.
\end{theorem}
\begin{proof}
Choose $a$ with $\rho_*<a<R/2$. The open sets
$U_C=\{x:\dist(x,C)<a\}$ cover the plane. Their intersection graph
is connected. Otherwise unions over different graph components would
be disjoint nonempty open sets partitioning the connected plane. If
$U_C\cap U_D\ne\varnothing$, then $\ell(C,D)<2a<R$.
The unrestricted gap graph is therefore connected, and
Lemma~\ref{lem:obstruction-descent} proves connectivity of
$\mathcal G_R$.

Choose a root and a spanning tree of the finite quotient, and lift the
tree to choose one representative copy at each vertex. Orient edges
arbitrarily and lift each from its chosen source copy. Write its target
as the chosen target representative translated by $d_e\in\Lambda$;
tree translations vanish.
Let $\Gamma$ be the integer group generated by the non-tree $d_e$.
For any $\lambda\in\Lambda$, connectivity of the lifted graph gives
a finite path from the root copy to its translate by $\lambda$.
Its projection is a closed walk. Summing its signed edge translations
shows $\lambda\in\Gamma$. Conversely each $d_e$ is a period, so
$\Gamma\subset\Lambda$. Hence $\Gamma=\Lambda$.
This is the elementary path-lifting argument for a periodic graph;
its graph-cover interpretation is classical, as in \cite{Gross}.
\end{proof}

Each edge in this theorem is an admissible local normal channel for
smooth positively curved obstacles. The closest pair gives a positive
quadratic stationary minimum, and clearance provides facing patches
and a short three-impact collar. The finitely many quotient edges have
some common positive local bounds for each fixed table. Such bounds
are not asserted uniformly over unrestricted degenerating tables.

\begin{definition}[A complete range-$R$ catalogue]\label{def:catalogue}
For every edge orbit of $\mathcal G_R$, record the two endpoint density
functions of a strictly active local return experiment. To avoid any
choice of a source, the catalogue may contain both return orientations;
one orientation for every edge also suffices. The two times, the local
origin, itinerary, and coherent reversal remain marks within each
record. Between records there are no obstacle identities, incident-vertex
labels, relative signs, contact offsets, or integer copy labels.
Duplicating an edge by recording its reverse changes neither connectivity
nor the cycle group. A full catalogue is supplied up to permutation of
records. A proper subcollection is not called complete.
\end{definition}

A known upper bound $\rho_0\ge\rho_*$ and any $R>2\rho_0$ specify
a sufficient finite acquisition range without knowledge of a period
basis or a desired spanning tree. The definition does not supply an
algorithm for locating and certifying every channel from partial data.
Its advantage over a designed network is that coverage and generation
are geometric consequences of one uniform range condition. Completeness
of the supplied catalogue and the covering bound remain explicit inputs.

\subsection{Saturation by all cycles, rather than one pair}
We give a finite arithmetic certificate for the group generated by all
observed cycles. All geometric vectors are columns. The matrix $B$ in
\eqref{eq:hnf-main} instead lists a basis of the dual integer lattice by
rows: its row lattice is $B^t\Z^2$, whose dual is $B^{-1}\Z^2$.
An integer matrix $H$ below is a column basis for a direct lattice.
These conventions are fixed throughout. Hermite and Smith normal forms
and their determinantal interpretation are standard; see \cite{Cohen},
Chapter~2. The rank-two index calculation needed here is proved below.

\begin{theorem}[The all-cycle index]\label{thm:all-cycle-index}
In the setting of Theorem~\ref{thm:descent-main}, let
$d_1,\ldots,d_m$ be all non-tree displacements, with
$D=[d_1\ d_2]$ invertible, $V=\covol\Lambda$, and
$n=|\det D|/V$. Put
\begin{equation}\label{eq:integer-columns}
 b_j=nD^{-1}d_j\in\Z^2,\qquad
 B_0=[b_1\ \cdots\ b_m],\qquad
 M=B_0\Z^m\subset\Z^2.
\end{equation}
Let $H$ be any integer column basis for $M$, and put
$g=\gcd_{i<j}|\det(b_i,b_j)|$, including zero minors in the gcd.
Then the recovered cycle lattice and its calibrated index are
\begin{equation}\label{eq:all-cycle-index}
 \Gamma=\sum_j\Z d_j=\frac1nDH\Z^2,\qquad
 q=[\Lambda:\Gamma]=\frac gn
    =\gcd_{i<j}\frac{|\det(d_i,d_j)|}{V}.
\end{equation}
All entries in the last gcd are nonnegative integers and $q\ge1$.
The data admit at most $\sigma_1(q)$ compatible lattices, obtained by
applying Lemma~\ref{lem:overlattice} to the basis $DH/n$ and volume $V$.
If $q=1$, the whole table is uniquely reconstructed even when no pair
of observed cycles is a period basis.
\end{theorem}
\begin{proof}
The quotient $\Lambda/D\Z^2$ has order $n$, so $n\Lambda\subset
D\Z^2$, proving integrality of all $b_j$. The first two columns are
$n(1,0)^t,n(0,1)^t$, so $M$ has rank two. Its column basis $H$ gives
the first identity in \eqref{eq:all-cycle-index}.

For completeness, write $B_0=HE$ with $E$ integral. Its columns
generate $\Z^2$, by the definition of $M$, so there is an integral
matrix $F$ with $EF=I_2$. The two-dimensional Cauchy--Binet formula
expresses $1$ as an integer combination of the two-column minors of $E$.
Their gcd is thus one. Since every two-column minor of $B_0$ equals
$\det H$ times the corresponding minor of $E$, it follows that
$g=|\det H|=[\Z^2:M]$. Existence of a basis of the integer subgroup
can equally be obtained by transposing the elementary Hermite reduction
in Lemma~\ref{lem:overlattice}.

Taking covolumes gives
$\covol\Gamma=|\det D|g/n^2=Vg/n$. The subgroup inclusion
$\Gamma\subset\Lambda$ makes this ratio the positive integer $q$.
Each pair $d_i,d_j$ consists of periods, so its determinant divided by
$V$ is integral; transformation by $D/n$ shows that their gcd is $g/n$.
Every compatible period lattice contains all of $\Gamma$ and has
volume $V$. The classification with this smaller index gives exactly
$\sigma_1(q)$ algebraic candidates before the physical tests. For
$q=1$ the only candidate is $\Gamma$.
\end{proof}

The cycle group is independent of the tree. Changing representative
copies adds a vertex translation difference to each edge; these cancel
on closed walks. Closed walks are integral combinations of fundamental
cycles, so changing the spanning tree changes generators, not $\Gamma$.
The certificate uses all cycles and does not require any of its many
possible generating sets to contain a two-element basis.

\begin{proposition}[Primitive generation without a primitive pair]
\label{prop:no-primitive-pair}
There are open analytic physical families with a connected covering
three-record collection for which $\Gamma=\Lambda$, but the three
pair determinants have normalized absolute values $2,3,5$.
\end{proposition}
\begin{proof}
Take one sufficiently small asymmetric analytic strictly convex body
$C$ and its translates by $\Lambda=\Z^2$. Observe the three nearest
normal bridges to copies at
\[
        d_1=(1,0)^t,\qquad d_2=(1,2)^t,\qquad d_3=(-1,3)^t.
\]
Each vector is primitive in $\Z^2$, so its open center segment contains
no lattice point. There are finitely many such segments modulo periods;
each has positive distance from all other lattice centers. For a
sufficiently small body the shortest body-to-body segments remain near
these center segments and are clear. The support functions used in
Proposition~\ref{prop:primitive} supply asymmetric analytic examples.
The quotient graph has one vertex and three loops. The determinants are
$2,3,5$, and their gcd is one. Directly,
$d_3-d_2+2d_1=(0,1)^t$, while $d_1=(1,0)^t$.
Thus all cycles generate $\Z^2$ but no observed pair does.
The three gaps, clearances and asymmetric shape persist under small
analytic perturbations of $C$, and the displacement vectors remain the
three period translations. This proves the open-family assertion.
\end{proof}

\subsection{Whole-table rigidity of the complete catalogue}
\begin{theorem}[Canonical-catalogue rigidity]\label{thm:catalogue-rigidity}
Let an analytic periodic dispersing table have asymmetric, pairwise
noncongruent obstacle representatives. Let $\rho_0$ bound its covering
radius and let $R>2\rho_0$. Its complete range-$R$ catalogue of two
endpoint density functions per return record determines the entire
periodic table and its full translational period group, up to one
Euclidean isometry and permutation of representatives. No obstacle
identifications, incidence graph, contact registration, relative record
signs, period basis, integer copy labels, or primitive cycle pair are
supplied. Coverage and generation follow from completeness at this
range; the number of unvisited obstacle orbits is zero by the geometric
theorem, not by a test on omitted data.
\end{theorem}
\begin{proof}
Theorem~\ref{thm:local-main} reconstructs each smooth pair of arcs and
its area normalization. Analytic continuation and
Lemma~\ref{lem:incidence} reconstruct complete images, their incidence
and a tree placement, as soon as connected coverage is known.
Theorem~\ref{thm:geometric-completeness} supplies exactly this coverage
and gives $\Gamma=\Lambda$. Consequently the calibrated index $q$ in
Theorem~\ref{thm:all-cycle-index} is one, and its finite construction
recovers a period basis $DH/n$. The local signs have cancelled through
unique shape matching, so only one ambient isometry remains. On the
shape-separated class the lattice is the full translational period
group, by the remark after Theorem~\ref{thm:descent-main}.
\end{proof}

Analyticity, asymmetry and distinct shapes enter the inverse step, not
the geometric completeness theorem. The latter holds for disks and
congruent representatives as well, but shape equality alone cannot
resolve their multiplicities or orientations. The retained registered
theorems address those different information categories. Likewise the
prime-index ambiguity in Proposition~\ref{prop:ambiguity} remains an
exact physical obstruction for its two selected records. Such a list
cannot be a complete catalogue at a range exceeding $2\rho_*$, because
its lifted cycle group has index greater than one. Additional short
clear channels must exist; the obstruction-descent theorem proves this
without being told their integer labels.

\subsection{Finite observations and stable saturation}
The following extension uses the same physical comparison method and
histogram estimator as Section~\ref{sec:stability}. Replace its assumption
of a primitive cycle pair by $\Gamma=\Lambda$. Keep the compact analytic,
shape and symmetry margins, the bounded lattice and inverse, positive
covolume, local physical bounds, bounded number $J$ of records, and onset
brackets. One may take compact classes of complete range-$R$ catalogues
with $\rho_*\le\rho_0<R/2$. For continuous statistical parametrizations,
fix a positive gap-to-cutoff margin for pair types near $R$, positive
clearance for the recorded branches, and charts in which the finite
catalogue is indexed before quotienting by permutations and reversals.
Enlarge the fixed presentation bound in \eqref{eq:loss} to include the
bases $DH/n$ arising from the bounded integer arrays below; their inverse
norms are bounded by the positive covolume bound. Thus the same loss
compares actual images and unmarked lattices, not independent edge frames.
These are prior regularity conditions; they do not confer a method for
certifying catalogue completeness from samples.

\begin{lemma}[Stable all-cycle saturation]\label{lem:stable-saturation}
For two compatible physical candidates in this bounded class, sufficiently
small errors in corresponding images and cycle vectors give the same
integers $n$ and columns $b_j$ in \eqref{eq:integer-columns}. Their
all-cycle indices $q$ therefore agree. If $q=1$, they have corresponding
period bases $DH/n$ and $\widetilde D H/n$, with the same integer $H$,
and the difference of these bases is at most a class-dependent constant
times $\|D-\widetilde D\|$.
\end{lemma}
\begin{proof}
Choose any independent pair of true cycle vectors. Its determinant has
absolute value at least the positive covolume lower bound, since it is
an integer multiple of that covolume. The vectors are uniformly bounded:
tree paths have bounded length and every record has bounded geometry.
Thus $D^{-1}$ and the positive integer $n$ are uniformly bounded.
Lemma~\ref{lem:integer-lock} gives equality of $n$ and every $b_j$ once
the physical discrepancy is small enough. The same deterministic
integer basis algorithm then chooses the same $H$ for both candidates.
There are only finitely many bounded integer arrays of at most $J$
columns, so the norms of the resulting $H$ have a uniform bound.
Theorem~\ref{thm:all-cycle-index} gives equality of $q$ and, if $q=1$,
the two period bases. Their difference is
$(D-\widetilde D)H/n$, which has the stated bound. This argument does
not form an integer group from raw noisy real vectors.
\end{proof}

\begin{theorem}[Finite-histogram recovery without a primitive pair]
\label{thm:catalogue-stability}
On the compact classes just specified, the experiment of
Theorem~\ref{thm:stat-main} reconstructs the incidence graph and the
whole table, including its unmarked period lattice, without a primitive
cycle pair. With at most $CJN$ independent phase preparations, including
failures, its error with probability at least $1-\delta$ is at most
\begin{equation}\label{eq:catalogue-stat-rate}
 C\left(\frac{\log(CJN/\delta)}{N}\right)^{\beta\theta/(2\beta+6)},
       \qquad 0<\beta\le1,\quad 0<\theta<1.
\end{equation}
The constants, threshold in $N$, and continuation exponent depend on the
entire prior class. The statement applies in particular to regular
complete-catalogue classes at range $R>2\rho_0$. It is a conditional
regularization theorem, not a minimax or computational-efficiency claim.
\end{theorem}
\begin{proof}
Use the retained pilot and the two growing endpoint histograms per
record. The concentration and cell reconstruction proof in
Section~\ref{sec:stability} gives simultaneous density error
\[
 e_N\le C\{h^\beta+\sqrt{r_N}h^{-3}+r_Nh^{-4}\},\qquad
 r_N=\log(CJN/\delta)/N.
\]
Its admissible approximate minimizer and Borel selection require no
primitive pair. The same physical prior charts and bounded preparations
justify continuity and conditional concentration at the adaptive times.
Lemma~\ref{lem:sync-stable} supplies an incidence isomorphism, a common
tree gauge, and errors at most $Ce_N^\theta$ in the complete images,
cycle vectors and volume. Apply Lemma~\ref{lem:stable-saturation} in
place of the old index-one-pair step. The true all-cycle index is one,
so the candidate has the same index and corresponding period basis.
The table loss is at most $Ce_N^\theta$. Taking
$h=r_N^{1/(2\beta+6)}$ proves \eqref{eq:catalogue-stat-rate} with the
same preparation count. The finite alignment and integer-basis choices
are part of the proof; the estimator is not supplied a correct incidence
partition, cycle pair or lattice basis.
\end{proof}

For an arbitrary partial but covering collection the same argument locks
its positive all-cycle index and its finite algebraic period candidates;
uniqueness is not asserted when physical alternatives remain. Exact
congruence, continuation and admissible minimization remain operations
on function-valued data or a compact prior. The present extension removes
a network-generation hypothesis for a complete bounded-range protocol;
it does not turn that protocol into an ordinary collision-count sensor
or a marked length spectrum.
